\documentclass[accentcolor=tud4b,colorbacktitle,inverttitle,landscape,german,presentation,t]{tudbeamer}
\mathcode`\,="013B%L�scht das Leerzeichen nach dem Komma in Dezimalzahlen
%\usepackage{multimedia}

%\usepackage{media9}

%\usepackage{movie15}


\usepackage{hyperref}
\usepackage{bbding}
\usepackage{units}
\usepackage{epstopdf}
\usepackage{subfigure}
\usepackage{caption}
\usepackage{ngerman}
\usepackage{color}
\usepackage{rotating}
\usepackage{amsmath}
\usepackage[latin1]{inputenc}   % deutsche Umlaute
\usepackage[T1]{fontenc}        % Darstellung von Schrift als echte T1 Fonts und nicht als Bilder
\usepackage{graphicx}
\graphicspath{{bilder/}}
% \usepackage{multimedia}
% \usepackage[square]{natbib}
% \usepackage{bibgerm}
\usepackage{tikz}

\usetikzlibrary{shapes.geometric,shapes.arrows,decorations.pathmorphing}
\usetikzlibrary{matrix,chains,scopes,positioning,arrows,fit}
\usetikzlibrary{calc}
\usetikzlibrary{patterns}

% \usetikzlibrary{patterns}
 \usetikzlibrary{matrix}
 \usetikzlibrary{decorations.pathmorphing}
 \usetikzlibrary{arrows,shapes.geometric} % fadings, ???
 \usepackage{fancybox}
 \usepackage{amsmath}            % Mathematik-Umgebungen
\usepackage{bm}                 % fette Mathe Buchstaben
\usepackage{amssymb}            % spezielle Mathe-Symbole


\usepackage{color}

%
% Absolute positioning
% beamer paper size: 12.80cm x 9.60cm
%
\usepackage[overlay,absolute]{textpos}
\setlength{\TPHorizModule}{10mm}
\setlength{\TPVertModule}{\TPHorizModule}
\textblockorigin{5mm}{25.5mm}
\setlength{\parindent}{0pt}

\newcommand{\leftbox}[1]{
 \begin{textblock}{5.7}(0, 0)
   #1
 \end{textblock}
}

\newcommand{\rightbox}[1]{
 \begin{textblock}{5.7}(6.1, 0)
   #1
 \end{textblock}
}

\newcommand{\fullbox}[1]{
 \begin{textblock}{11.8}(0, 0)
   #1
 \end{textblock}
}

% \usepackage{calc}
\newenvironment{arbbox}[2]
{%\setlength{\boxwidth}{11.8cm minus \real{#1}}
\begin{textblock}{11.8cm}(#1, #2)}
{\end{textblock}}




%
% Macros
%
% \renewcommand{\vec}[1]{\ensuremath{\bm{#1}}}
\renewcommand{\vec}[1]{\ensuremath{\boldsymbol{#1}}}

\begin{document}
\title[]{Titel}
\subtitle{Namen}

\author[Gruppe X]{Gruppe X}
\institute[Regelungsmethoden und Robotik | Projektseminar Robotik]{Regelungsmethoden und Robotik | Projektseminar Robotik}




% \logo{\color{tudtextaccent}TU Darmstadt-}
\logo{\includegraphics{rmr2}}

\date{15.07.2017}
%Titelfolie%%%%%%%%%%%%%%%%%%%%%%%%%%%%%%%%%%%%%%%%%%%%%%%%%%%%%%%%%%%%%%%%%%%%%%%%%%%%%%%%%%%%%%%%%%%%%%%%%%%%%%%%%%%%%%%%%%%%%%%%%%%%%%%%%%%%%%%%%%%%%%%%%%%
\begin{titleframe}





\end{titleframe}

%Hauptfolien%%%%%%%%%%%%%%%%%%%%%%%%%%%%%%%%%%%%%%%%%%%%%%%%%%%%%%%%%%%%%%%%%%%%%%%%%%%%%%%%%%%%%%%%%%%%%%%%%%%%%%%%%%%%%%%%%%%%%%%%%%%%%%%%%%%%%%%%%%%%%%%%%%%
%%%%%%%%%%
\begin{frame}{\\Distortion example}

 \vspace{-1cm}
 \begin{align*}
 x &=x_d+ (x_d-x_c)(1+K_1r^2+K_2r^4)+P_1\left( r^2+2(x_d-x_c)^2\right) +2P_2(x_d-x_c)(y_d-y_c)\\
 y &= y_d+(y_d-y_c)(1+K_1r^2+K_2r^4)+2P_1(x_d-x_c)(y_d-y_c)+P_2\left(r^2+2((y_d-y_c)^2\right)
 \end{align*}


 \leftbox{
 \vspace{2cm}
 \textbf{Radial distortion:}
 \begin{align*}
 K_n = n^{th} \text{ radial distortion coefficient}
 \end{align*}
}
\rightbox{
 \vspace{2cm}
 \textbf{Tangential distortion:}
  \begin{align*}
 P_n = n^{th} \text{ tangential distortion coefficient}
 \end{align*}

}
\vspace{2cm}
\begin{align*}
(x_d,~y_d) &= \text{ distorted imaage point as projected on image plane}\\
(x,~y) &= \text{ undistorted imaage point as projected on image plane}\\
(x_c,~y_c) &= \text{ distortion center}\\
r &= \sqrt{(x_d-x_c)^2+(y_d-y_c)^2}
\end{align*}
\end{frame}

\begin{frame}{\\Distortion example}

 \vspace{-1cm}
 \begin{align*}
 x &=x_d+ (x_d-x_c)(1+K_1r^2+K_2r^4)+P_1\left( r^2+2(x_d-x_c)^2\right) +2P_2(x_d-x_c)(y_d-y_c)\\
 y &= y_d+(y_d-y_c)(1+K_1r^2+K_2r^4)+2P_1(x_d-x_c)(y_d-y_c)+P_2\left(r^2+2((y_d-y_c)^2\right)
 \end{align*}


 \leftbox{
 \vspace{2cm}
 \textbf{Radial distortion:}
 \begin{figure}
 \includegraphics[scale=0.3]{distortions_inline}
 \caption{radial distortions}%source
 \end{figure}
}
\rightbox{
 \vspace{2cm}
\textbf{Tangential distortion:}
 \begin{figure}
 \includegraphics[scale=0.275]{tangential_dist}
 \caption{first order tangential distortion }%source
 \end{figure}

}

\end{frame}



\begin{frame}{\\pixel size detection}

%\vspace{4cm}
\leftbox{
 %\vspace{2cm}
 \textbf{pixelSize = 1:}
 \begin{figure}
 \includegraphics[scale=0.1]{pixelSizeImg1}
  \end{figure}
}
\rightbox{
 %\vspace{2cm}
\textbf{pixelSize = 8:}
 \begin{figure}
 \includegraphics[scale=0.1]{pixelSizeImg2}
 \end{figure}

}
 
\end{frame}

\begin{frame}{\\Center point estimation}

 \begin{figure}
 \vspace{-0.5cm}
 \includegraphics[scale=0.15]{pixelSizeImg2_FOV}
 \end{figure}
 \leftbox{
  \vspace{4.5cm}
  \begin{align*}
  \textbf{x}_c =\frac{\sum\limits_{k=1}^{n} \textbf{x}_k}{n}
  \end{align*}
 }
 \rightbox{
  \vspace{5cm}
  %\begin{align*}
  $n$~~~:~number of seen pixels\\
  $\textbf{x}_k$~:~position of seen pixel
  %\end{align*}
}
\end{frame}


\begin{frame}{\\Results}
\leftbox{
 %\vspace{2cm}
 \textbf{Ground truth:}
 \begin{figure}
 \includegraphics[scale=0.2]{gt}
 \caption{all white lines that we're drawn and seen on the screen}
  \end{figure}
}
\rightbox{
 %\vspace{2cm}
\textbf{Mapped Image: }
 \begin{figure}
 \includegraphics[scale=0.2]{mappedImage}
 \caption{the seen lines after they were mapped by the algorithm}
 \end{figure}
 }
\end{frame}

\begin{frame}{\\Comparison\\ \small{Substraction of ground truth and mapped image}}
\begin{figure}
\vspace{-0.5cm}
\includegraphics[scale=0.4]{substracted}
\caption{difference of both images\\ $824$ of $290,191$ pixels do not fit}
\end{figure}

\end{frame}






\end{document}
%% High-resolution temporal profilometry using Fourier estimation (PDF Download Available). Available from: https://www.researchgate.net/publication/236330080_High-resolution_temporal_profilometry_using_Fourier_estimation [accessed Jul 16, 2017] Source image tangential dist